\documentclass[12pt]{article}
\usepackage[utf8]{inputenc}
\usepackage[T1]{fontenc}
\usepackage{hyperref}
\title{Waur\'a Language Resources for K-12 Curriculum Development}
\author{Compiled by Claude Sonnet 5 (claude-sonnet-5) for Andr\'as Kornai}
\date{}
\begin{document}
\maketitle

\section{Introduction}
Waur\'a, or Wauja (ISO 639-3 \texttt{wau}, Glottocode waur1244), is an Arawakan language spoken in the Xingu Indigenous Park, Mato Grosso, Brazil. Estimates run from about 320 to 410 speakers. The search found no dictionary, textbooks or school materials, so this report is brief.

\section{Primary Educational Resources}
None were found.

\section{Linguistic and Reference Materials}
\begin{itemize}
\item Wikipedia, Waur\'a language: \url{https://en.wikipedia.org/wiki/Waur%C3%A1_language}
\item Wikipedia, Wauja (the people): \url{https://en.wikipedia.org/wiki/Wauja}
\item Endangered Languages Project entry: \url{https://endangeredlanguages.com/node/38101}
\item Journal article in \emph{Estudos Lingu\'isticos}, found in search; I did not review its contents: \url{https://revistadogel.emnuvens.com.br/estudos-linguisticos/article/download/1226/779/3333}
\item Journal article in \emph{Estudos da Linguagem}, found in search; I did not review its contents: \url{https://periodicos2.uesb.br/index.php/estudosdalinguagem/article/download/1013/864/1699}
\end{itemize}
Wikipedia cites Postigo's 2014 work, \emph{L\'ingua Wauja (Arawak): uma descri\c{c}\~ao fonol\'ogica e morfossint\'atica}. I did not find a link to it in the search results.

\section{Religious and Community Resources}
None were found.

\section{Digital and Audio Resources}
None were found.

\section{Curriculum Development Context}
Waur\'a is not known to be taught in schools. Brazilian indigenous school education in the Xingu is a possible avenue, but no evidence was found in this search. Contact with Wauja communities and Brazilian linguists would be needed.

\end{document}